\begin{lemma}[one common signed cap convention]\label{push:src:C.3.22}
Retain the relative Pontryagin number
$p_{\mathrm{rel}}=-24$ of \ref{push:src:C.3.18} and the actual relative
Euler square $s_\delta$ of \ref{push:src:C.3.20}. With the single
cut-and-reglue orientation sign $\epsilon_{\rm cap}$ used
in \ref{push:src:C.3.20}, the two cap pairings satisfy
\[
 \langle p_1(M_\psi)x,[M_\psi]\rangle
       =\epsilon_{\rm cap}\,p_{\mathrm{rel}},
 \qquad
 k_\psi=\epsilon_{\rm cap}\,3s_\delta.
\]
In particular
\[
 J_\psi=\epsilon_{\rm cap}
              \frac{p_{\mathrm{rel}}-12s_\delta}{24}.
\]
This lemma does not assign a value to $s_\delta$.
\end{lemma}
\Uses{the hypotheses, notation and definitions stated here; \cref{push:src:C.3.18}; \cref{push:src:C.3.20}; \cref{push:src:C.3.19}}

\begin{proof}
We fix the normalization by connections, so that a boundary
frame and a coordinate clutch are not silently identified.
For a real orthogonal connection with curvature $F$, use
\[
 p_1(F)=-\frac{\operatorname{tr}(F\wedge F)}{8\pi^2},
 \qquad
 \operatorname{CS}_{p_1}(\omega)
 =-\frac{\operatorname{tr}(\omega\wedge d\omega
                          +\tfrac23\omega^3)}{8\pi^2}.
\]
For $K$ a change of coordinates,
\[
 \operatorname{CS}_{p_1}(K^{-1}dK)
 =\frac{\operatorname{tr}(K^{-1}dK)^3}{24\pi^2}
 =:\tau_{p_1}(K).
\]
This is the suspension used for a mapping torus oriented
with its interval parameter first. Indeed interpolating
the connection from zero to $K^{-1}dK$ gives curvature
$dt\wedge\omega+(t^2-t)\omega^2$, whose first Pontryagin
form integrates in $t$ to the displayed three-form.

In contrast, let a boundary frame be
$\mathfrak f=\mathfrak t B$ relative to a global stable
frame $\mathfrak t$. The connection for which
$\mathfrak f$ is parallel has matrix $-dB B^{-1}$ in
the $\mathfrak t$ coordinates. Therefore
\[
 p_1(E,\mathfrak f)[M,\partial M]
 =-\int_{\partial M}\tau_{p_1}(B).
\]
Equivalently, capping the boundary frame by a standard
frame uses the inverse coordinate clutch $B^{-1}$.
This minus sign is forced by the connection computation;
it is not an orientation chosen to obtain a cancellation.

Now use exactly the change-of-frame convention of C.3.19:
\[
 \mathfrak t=\mathfrak f_{\mathrm{aff}} A,
 \qquad
 G=A(\psi p)^{-1}A(p)
\]
after contracting the affine derivative shear. Set
\[
 a_0=\int_{T_{\gamma,u,w}}\tau_{p_1}(A).
\]
Here $B=A^{-1}$ in the boundary-frame formula. Since
the boundary orientation of $M$ is
$\partial M=-T_{\gamma,u,w}$, that formula gives
\[
 p_{\mathrm{rel}}
 =-\int_{\partial M}\tau_{p_1}(A^{-1})
 =-a_0.
\]

For the actual annulus use the orbit orientation
$dg\wedge du\wedge dw\wedge dr\wedge da$ and the
Euler curvature $E=dg\wedge da$, with $\chi$ decreasing
from one to zero as $r$ increases. The relative trace
calculation of \ref{push:src:C.3.19} has representative
\[
 \tau_{p_1}(G)
 =-d\chi\wedge\beta
 \quad\hbox{in relative cohomology},
 \qquad
 \int_{T^2_{u,w}}\beta=a_0.
\]
Consequently
\[
 \int_P\tau_{p_1}(G)\wedge E=-a_0=p_{\mathrm{rel}}.
\]
For example, in the only contributing term the form is
$-a_0\,d\chi\wedge du\wedge dw\wedge dg\wedge da$;
relative to the displayed orbit orientation its coefficient
is $a_0\chi'(r)$, and the integral of $\chi'$ is $-1$.

This is also the Pontryagin--Euler pairing of the actual
mapping torus in \ref{push:src:C.3.20}. Its line curvature restricts to
$E+t\Delta$ on the $t$-slices. The possible extra term
has integral zero, because
\[
 \int_{DP}\tau_{p_1}(G)\wedge\Delta
 =\int_{DP}\tau_{p_1}(G)\wedge d\eta=0.
\]
Thus no line-clutch correction changes the preceding
pairing. The same cut-and-reglue bordism of \ref{push:src:C.3.20} is a
smooth bordism of tangent bundles and of the specified
line bundle. Both characteristic polynomials $p_1x$
and $x^3$ therefore acquire the same boundary-orientation
sign $\epsilon_{\rm cap}$. Their identity-cap and
identity-mapping-torus values vanish. Combining the
computed mapping-torus pairing with the cubic calculation
of \ref{push:src:C.3.20} gives the two asserted signed equations and
then their stated combination.
\end{proof}

